\documentclass[11pt]{article}
\usepackage[margin=1in]{geometry}
\usepackage[T1]{fontenc}
\usepackage{lmodern}
\usepackage{amsmath,amssymb}
\usepackage{booktabs}
\usepackage{enumitem}
\usepackage[round]{natbib}
\usepackage[hidelinks]{hyperref}
\usepackage{microtype}
\usepackage{xcolor}
\newcommand{\todo}[1]{\textcolor{red}{[TODO: #1]}}
\newcommand{\bp}{\,bp}

\title{When Realistic Market Simulators Give Unreliable Execution Advice:\\[4pt]
\large Sampling Error, Misspecification, and Decision Regret}
\author{Gaurav Chattree\thanks{Experiments, code and manuscript drafting were assisted by an AI coding agent (Claude Code). Every number in this draft is produced by scripts in the accompanying repository (branch \texttt{simulator-decision-study}). \todo{co-authors, affiliations, employer approval}}}
\date{Draft, 16 November 2026}

\begin{document}
\maketitle

\begin{abstract}
Limit-order-book simulators are increasingly used to choose how to execute large orders, and they are increasingly
validated by how realistic they look. We ask whether realism, or any other evidence available before trading, tells a
trader when a simulator's execution advice can be trusted. On Nasdaq data for three stocks, simulators of similar
realism recommend opposite strategies, a disagreement of up to 7.6\bp{} replicated under pre-registration. In 26
synthetic reference markets with known costs, no observational realism score from LOB-Bench or LOB-Impact predicts good
advice, and optimising over a fitted simulator whose mechanisms differ from the market's is worse than always trading
evenly (TWAP). Decomposing regret shows the excess is structural rather than sampling: removing a single mechanism
(fundamental anchoring) from the true generator makes its advice worse as the simulation budget grows, from 7.9 to
11.2\bp{} between 10 and 300 runs. A pre-registered second study then tests remedies on 48 untouched markets: trust rules
using only observable evidence are significantly worse than TWAP (+1.00\bp{} [+0.24, +1.51]); a simulator with the
missing mechanisms added improves on uncorrected baselines on development markets but not significantly on the test
markets; a learned diffusion generator (TRADES) and a direct feature-to-decision estimator do not help either. The
decision-relevant mechanisms are, however, identifiable from observable data within a model family: the obstacle is
the candidate simulators' structure, not the information in the data. We release the reference markets, protocol and
code as a benchmark for decision reliability.
\end{abstract}

\section{Introduction}

Market simulators are now routinely proposed as testbeds for execution and trading algorithms, from agent-based
platforms such as ABIDES \citep{byrd2019abides,amrouni2021abidesgym} to point-process models \citep{bacry2015hawkes},
queue-reactive models \citep{huang2015queue} and learned generators \citep{coletta2021towards,berti2025trades}. Their
validation has largely been observational: a simulator is judged by how closely the order flow and book it generates
match real data \citep{vyetrenko2020getreal,nagy2025lobbench}. Recent work has added interventional checks, notably the
response of a generator to injected metaorders \citep{nigmatulin2026lobimpact}, and has begun to measure the value of
simulated experimentation by decision regret \citep{lu2026crashgym}.

A trader who uses a simulator to choose an execution schedule needs something narrower: confidence that the strategy
the simulator recommends is better than a simple default, in this market. We study that question directly, in three
parts.
\begin{enumerate}[leftmargin=*]
\item \textbf{Agreement.} Do simulators that look equally realistic on real data agree on the execution decision?
\item \textbf{Reliability.} In markets where the true cost of every strategy is known, which properties of a simulator
predict good advice, and is its error sampling noise or structure?
\item \textbf{Validation.} Can evidence available before trading tell when to trust a simulator's advice, and does
correcting its missing mechanisms make the advice reliable?
\end{enumerate}

\paragraph{Contributions.}
(i) A pre-registered demonstration on Nasdaq data that similarly realistic simulators recommend opposite strategies,
with a mechanism-level explanation (Section~\ref{sec:goal1}). (ii) A reference-market benchmark with known costs in
which we show that neither observational realism nor LOB-Impact's free-run scores predict execution regret, that
selecting the cheapest predicted strategy amplifies model error, and that the excess regret is structural
(Section~\ref{sec:goal2}). (iii) A decisive single-mechanism experiment: the true generator with one mechanism removed
becomes more confidently wrong with more simulation. (iv) A pre-registered test on 48 untouched markets of three
remedies, trust rules, mechanism correction and a learned generator, all of which fail to beat TWAP, together with the
finding that the decision-relevant mechanisms are nonetheless identifiable from observable data
(Section~\ref{sec:trust}).

\section{Related work}
\label{sec:related}

\paragraph{Simulator validation.} LOB-Bench \citep{nagy2025lobbench} standardises observational realism scores for
generative order-book models; LOB-Impact \citep{nigmatulin2026lobimpact} measures emergent impact under injected
metaorders. We use both as candidate validation measures and test whether they predict decision quality.
\paragraph{Decision-focused evaluation.} CrashGym \citep{lu2026crashgym} scores LLM agents' crash interventions by
regret against counterfactual outcomes on replayable synthetic markets, with a laboratory that is a copy of the true
world. Our candidates are fitted, misspecified simulators, the task is optimal execution, and the object of study is
the simulator and its validation rather than an agent's experimentation.
\paragraph{Causal and generative impact models.} Causal foundation models estimate impact under hidden confounding
\citep{thumm2026causal}; regime-conditioned diffusion generates counterfactual trajectories \citep{wang2026difflob};
self-calibrating simulators track recorded data \citep{zhong2026evomarket}; compound Hawkes models add realistic order
sizes and spread dependence \citep{jain2024compound}. We include compound Hawkes and TRADES as modern baselines.
\paragraph{Robust execution and decisions under ambiguity.} Robust execution under misspecified dynamics
\citep{schied2013robust}, decisions under partial identification \citep{manski2011choosing}, safe improvement over a
baseline policy \citep{laroche2019spibb,joshi2024safe} and policy-specific model selection \citep{dong2023local} are
established. We do not propose a new robust optimiser; we measure, in a setting with known answers, whether the
evidence such methods rely on is informative for execution. \todo{verify Joshi et al.\ 2024, Dong et al.\ 2023; add
partial-identification and confounded offline RL references from the novelty review after checking them}

\section{Setup}
\label{sec:setup}

\paragraph{Task and strategies.} Sell a fixed quantity (5--15\% of a ten-minute volume in the synthetic markets; 10\%
and 30\% on real data) over ten minutes in 20 decisions, using market orders. Twelve strategies span schedules (TWAP,
selling at once, front- and back-loaded, risk-averse), participation and liquidity-aware rules, and signal rules that
speed up or slow down with the last slice's price move and order-flow imbalance. Cost is implementation shortfall
against the arrival mid, paired with an unperturbed twin run on the same seed.
\paragraph{Regret.} For a market with true expected costs $C(\pi)$, a decision maker recommending $\hat\pi$ incurs
regret $C(\hat\pi)-\min_\pi C(\pi)$. True costs come from 300 reference seeds; where noted, the best strategy is
identified on seeds 0--149 and costs evaluated on 150--299 (cross-fitting).
\paragraph{Candidates.} On real data: Hawkes models (original, queue-aware, refill) and ABIDES (original and
calibrated), plus a TRADES case study. In synthetic markets: 29 candidates fitted to six 30-minute observable episodes
of each market, 28 Hawkes variants (spread closing, gap filling, multi-timescale kernels, impact scaling, stability
caps) and ABIDES calibrated by simulated moments; later, a queue-reactive candidate, compound Hawkes, a recalibrated
Hawkes (24 episodes), a mechanism-corrected Hawkes and TRADES.
\paragraph{Measures.} Observational: LOB-Bench order-book, price and impact scores; LOB-Impact free-run statistics
(sign memory, Hurst exponent, response ratio). Interventional: LOB-Impact build-up exponent, peak and relaxation;
the error in a simulator's own TWAP impact (``metaorder accuracy''); a decision-aware validation profile.
\paragraph{Inference and pre-registration.} The unit of replication is the market. Intervals come from a three-level
bootstrap (markets, designs shared across markets, seeds within market and design). Every confirmatory test was written
and checksummed (SHA-256) before its data existed; eleven design addenda and all deviations are listed in
Appendix~\ref{app:prereg}.

\section{Agreement: realistic simulators disagree on real data}
\label{sec:goal1}

On Nasdaq TotalView-ITCH data (INTC for discovery; CSCO primary and MSFT secondary confirmation), the original Hawkes
model recommends selling at once and the original ABIDES model recommends patience on all three stocks. The decision
statistic $D=C(\mathrm{TWAP})-C(\mathrm{immediate})$ is $+3.53$, $+2.20$, $+2.23$\bp{} for Hawkes and $-4.10$, $-4.95$,
$-2.36$\bp{} for ABIDES at 10\% size: pre-registered disagreement (opposite signs, gap at least 2\bp) held on all three
stocks (gaps 7.63, 7.15, 4.59\bp). The hope that improved versions converge did not replicate.

The disagreement has a measurable cause. The decision weighs the cost of selling at once (book depth) against the cost
of selling gradually (the trader's own impact). ABIDES-original prices selling at once 1.3--3.1 times too high and so
advises patience; Hawkes-original makes gradual selling move the price strongly (1.3--1.9 times impact measured with
synthetic metaorders) and so advises selling at once. Which simulator ``gets impact right'' depends on the benchmark:
against the literature square-root band, Hawkes-refill and ABIDES-original fall inside; against impact measured on
public trades, Hawkes-original is closest. Better order-book realism went with more extreme advice (Spearman $-0.90$,
$-1.00$, $-0.20$). Real data cannot say which advice is right, which motivates markets with known answers.

\section{Reliability: markets with known costs}
\label{sec:goal2}

\subsection{Design}
Twenty-six reference markets whose rules we wrote: queue-reactive markets (flow responds to queue sizes), agent-mix
markets (noise traders, fundamentalists anchored to a latent value, trend followers with a hidden regime, liquidity
providers who withdraw from one-sided flow), ten markets drawn at random for confirmation, and six well-specified
markets built from the candidates' own Hawkes family. A held-out latent-liquidity family (six markets, liquidity
revealed around a slowly relaxing reference price) was used only once, at the end of this stage.

\subsection{Realism does not predict good advice}
On the ten drawn markets (pre-registered, 30 seeds), no realism measure predicted low regret (C3 holds). Metaorder
accuracy correlated positively but did not meet its confirmatory test (C1: $+0.15$ [$-0.03$, $0.28$]); choosing by it
rather than by LOB-Bench order-book score lowered regret by $2.66$\bp{} [$-4.75$, $+0.03$] (C2). After a Holm
adjustment neither directional hypothesis is significant (one-sided $p=0.056$ and $0.029$). In the held-out family,
realism pointed the wrong way: the LOB-Bench order-book score correlated $-0.32$ [$-0.61$, $-0.02$] with good advice.

\subsection{Optimisation amplifies model error}
Selecting the cheapest predicted strategy picks the most flattering error: the chosen strategy's predicted gain over
TWAP is overstated by $10.4$\bp{} on average, against $1.6$\bp{} for a typical strategy, and only 23\% of non-TWAP
picks beat TWAP \citep{smith2006optimizer}. Simulators judged selling at once about as cheap as TWAP; it was 8\bp{}
worse and never optimal. Gating departures from TWAP on the predicted gain's uncertainty removes the damage but, tuned
on misspecified markets, never leaves TWAP.

\subsection{The excess regret is structural}
At 100 decision seeds, choosing by the candidates' ensemble was worse than TWAP in agent-mix markets by $6.12$\bp{}
[$3.33$, $8.96$] and in the 20 misspecified markets by $3.02$ [$1.31$, $5.19$]; in the well-specified markets it tended
to beat TWAP ($-2.43$ [$-5.66$, $+1.54$]). Comparing four decision makers (true generator or candidates, 30 or 100
seeds) attributes almost all excess regret in agent-mix markets to structure, $+6.72$\bp{} [$4.23$, $7.83$] ($+9.11$
cross-fitted), against $+0.41$ [$-0.39$, $1.37$] to sampling. Budget curves show the difference: the true generator's
regret falls with more seeds (4.01, 2.31, 1.90\bp{} at 5, 30, 100 seeds) while a single candidate's rises (7.07, 8.15,
8.62). The held-out family replicates every conclusion: TWAP 0.45\bp, true generator 0.06, candidates' ensemble
$+1.78$ [$0.18$, $3.79$] worse than TWAP, structural share $+1.95$ [$0.59$, $3.04$].

\subsection{One missing mechanism}
To isolate structure from fitting, we take the true generator of four agent-mix markets and switch off exactly one
mechanism. Without fundamental anchoring the generator's regret rises with the simulation budget (7.88, 8.73, 11.28,
11.19\bp{} at 10, 30, 100, 300 seeds); at 300 seeds it never selects the true best strategy and its cost ranking is
anti-correlated with the truth ($-0.14$); it overstates the gain from selling at once by 17.8\bp. Removing trend
followers or liquidity withdrawal hardly matters. The full true generator converges in three of the four markets; in
the fourth, heavy-tailed costs make even the truth's 300 seeds unable to separate the top strategies.

\subsection{Further checks}
Refitted to markets in which one mechanism is varied over five levels, the fixed designs picked the true optimum at
2--10\% of levels in agent-mix markets and met a pre-stated decision-boundary test in 88--100\% of cases. Over a
30-minute task with 60 decisions, more is at stake (TWAP leaves 1--2.3\bp) and the true generator captures most of it,
but no fitted design beats TWAP on average. The synthetic families match Nasdaq on diffusion and short-horizon return
dependence but show a far more persistent response to trades ($R(100)/R(1)$ of 9--17 against 1.9) and steeper impact
(exponent 0.64--1.04 against 0.33--0.43); conclusions describe markets with these mechanisms.

\section{Validation: can a trader tell when to trust a simulator?}
\label{sec:trust}

A second pre-registered study asks whether evidence available in real markets can tell when a simulator's advice is
safe, and whether correcting its missing mechanisms helps. Rules and models were developed on 42 development markets,
read once on 16 Test~A markets, and confirmed on 48 new markets drawn without any filter (Test~B; 16 per family; size
set by a power calculation for a 0.5\bp{} improvement). Test~B was analysed once, after a checksummed freeze.

\paragraph{Trust rules.} Candidate rules depart from TWAP only when the fixed designs agree, when their advice is stable
across fitted variants, when the predicted gain exceeds its uncertainty, or when a design reproduces the market's
mechanism diagnostics (variance ratios, response to trades, depth response to one-sided flow). One rule, agreement, was
frozen on the development markets. Even the true generator fails its own diagnostic on 20--47\% of samples, so the
diagnostics are too blunt to separate good simulators from bad ones.

\paragraph{Identifiability.} In an 81-market grid crossing anchoring, trend following, withdrawal and replenishment,
only 1.0\% of pairs of markets requiring different decisions were observationally indistinguishable from six episodes
(median classifier AUC 0.99); withdrawal was the hardest mechanism to see. The information needed for the decision is
present in observable data.

\paragraph{Mechanism correction.} A Hawkes candidate with fundamental anchoring and liquidity withdrawal added, its
extra parameters fitted by simulated moments, improved on the original design by $1.60$\bp{} [$-3.22$, $-0.05$] and on
compound Hawkes by $2.25$ [$-3.81$, $-0.55$] on the development markets, but remained worse than TWAP ($+2.43$ [$1.18$,
$3.72$]); recalibrating on four times more data did not help, and compound Hawkes was no better than the study's
designs. The corrected model's own parameters are not identified from six episodes: settings with near-identical
observable features predicted TWAP costs of 67 and 32\bp{} (true 7.2) and different best strategies.

\paragraph{Test B.} Table~\ref{tab:testb} reports the confirmatory results. No hypothesis is supported. The frozen rule
departed from TWAP in 15\% of cases and those departures cost money; the corrected model was better than both Hawkes
baselines in every family but not significantly after adjustment.

\begin{table}[t]
\centering
\caption{Test~B: 48 untouched markets, cross-fitted regret (bp). Hypotheses one-sided, market bootstrap with decision
seeds resampled, Holm at family-wise 0.025.}
\label{tab:testb}
\begin{tabular}{lrrrr}
\toprule
 & All 48 & Agent-mix & Latent & Queue-reactive \\
\midrule
No-simulator rule & $-0.38$ & $-1.10$ & $0.01$ & $-0.04$ \\
Always TWAP & $0.37$ & $1.01$ & $0.09$ & $0.00$ \\
Frozen trust rule & $1.37$ & $3.85$ & $0.26$ & $0.00$ \\
Corrected Hawkes & $2.16$ & $5.87$ & $0.59$ & $0.02$ \\
Original Hawkes & $2.98$ & $6.83$ & $1.77$ & $0.35$ \\
Compound Hawkes & $3.50$ & $8.73$ & $1.71$ & $0.07$ \\
\midrule
\multicolumn{5}{l}{H10.1 rule $-$ TWAP: $+1.00$ [$+0.24$, $+1.51$]; H10.2 rule $-$ no-simulator: $+1.75$ [$+0.70$, $+2.52$]} \\
\multicolumn{5}{l}{H10.3 corrected $-$ original: $-0.82$ [$-2.00$, $+0.51$], $p=0.153$; H10.4 corrected $-$ compound: $-1.34$ [$-2.61$, $+0.17$], $p=0.041$} \\
\bottomrule
\end{tabular}
\end{table}

\paragraph{A learned generator and a direct estimator.} TRADES, trained on each of three development markets with four
times the data of the other candidates, recommended selling at once in every market; its predicted gains over TWAP were
off by 30--34\bp, the least accurate of any candidate. A gradient-boosted estimator mapping observable features directly
to strategy cost differences, evaluated leaving one market family out, beat TWAP on no family with headroom and was
significantly worse than TWAP on the held-out queue-reactive family; the no-simulator rule, which reuses the strategy
that won on the training families, also failed out of family (6.80\bp{} against TWAP's 4.22 on agent-mix markets).

\section{Discussion}
\label{sec:discussion}

The evidence points one way. Realism, interventional fidelity scores, agreement between candidates, stability, and
mechanism diagnostics computed from observable data do not tell a trader when a misspecified simulator's execution
advice is reliable, and optimisation turns the simulator's errors into confident recommendations. Because the
decision-relevant mechanisms are identifiable from the same data, the obstacle is the candidate's structure: fitted
families have directions along which the fit to observable data stays good while the decision changes. Adding the
missing mechanism narrows the damage without making the simulator worth following. A practical reading: a simulator's
promised saving over TWAP is not evidence on its own, and a simulator should be trusted for execution only where its
mechanisms are known to match the market's.

\paragraph{Limitations.} Real-data results cover three large US stocks, eight public sample days and sell orders; real
data cannot say which simulator is right. The synthetic markets are of our own design and differ from Nasdaq in their
response to trades. Several comparisons rest on few markets; the TRADES comparison is descriptive (three markets, 20
seeds). Test~B used a reduced set of candidate designs, no own-execution or probe-trade rules, and no opportunity-rich
stress family; these and every other deviation were recorded before the data were read.

\section{Reproducibility}
All code, reference markets, pre-registrations with their checksums, and summary results are in the repository; raw
simulation outputs can be regenerated from their seeds. \todo{public release, licence, data availability statement
for the ITCH sample files}

\bibliographystyle{plainnat}
\bibliography{refs}

\appendix
\section{Pre-registration timeline and deviations}
\label{app:prereg}
\begin{itemize}[leftmargin=*]
\item 31 Oct: confirmatory replication on CSCO and MSFT (H1--H5). Goal~2 design.
\item 1--3 Nov: LOB-Impact design; addenda 1--5 (ABIDES rollouts, validation against the reference tool, designed and
drawn markets, confirmatory hypotheses C1--C4, decision rules, well-specified markets).
\item 5--7 Nov: addenda 6--9 (positive control, mechanism ablation, 100-seed reruns, held-out family, matched
mechanisms, 30-minute task, removed-mechanism experiment, budget curves).
\item 12--15 Nov: addendum 10 (trust rules, identifiability, mechanism correction, Test~B) with amendment 1 (diagnostic
band level); addendum 11 (TRADES); Test~B freeze.
\item 16 Nov: Test~B read once; exploratory direct estimator (gate fixed before running).
\end{itemize}
Deviations and corrected errors are listed in \texttt{simrel/reports/addendum10\_notes.md} and the study report.
\todo{condense the deviations table into the appendix}

\end{document}
